\section{Monodromy and the spin representation}\label{sec:monodromy}

The degree-zero part of the LLV algebra contains the rational
subalgebra
\[
 \mathfrak s\cong\mathfrak{so}(H^2(M;\QQ),q_M).
\]
Its representation on $\IH^\bullet(M;\QQ)$ integrates to the
simply connected spin group. Let $G$ be its image as a connected
rational algebraic group. The bracket formulas of Section
\ref{sec:algebra} give an isogeny
\[
 G\longrightarrow\SO(H^2(M;\QQ),q_M),\qquad \sigma\longmapsto\bar\sigma,
\]
such that
\begin{equation}\label{eq:spin-conjugacy}
 \sigma L_x\sigma^{-1}=L_{\bar\sigma x},\qquad
 \sigma\Lambda_x\sigma^{-1}=\Lambda_{\bar\sigma x},\qquad
 \sigma h=h\sigma.
\end{equation}
For a nondegenerate complex subspace $W\subset H^2(M;\CC)$,
write $G_W$ for the connected subgroup of $G_{\CC}$ integrating
$\mathfrak{so}(W)$, acting trivially on $W^\perp$ in $H^2$.

\begin{proposition}\label{prop:spin}
Suppose a group $\Gamma$ acts on $\IH^\bullet(M;\QQ)$ by graded
automorphisms preserving the intersection pairing and a lattice.
Suppose also that it acts on $H^2(M;\QQ)$ by isometries preserving
a lattice, and that these actions satisfy
$\gamma L_x\gamma^{-1}=L_{\gamma x}$. A subgroup of finite
index of $\Gamma$ acts on $\IH$ through $G(\QQ)$, compatibly
with its action on $H^2$. If the complex Zariski closure of its
action on $H^2$ contains $\SO(W)$, the closure of its action on
$\IH$ contains $G_W$.
\end{proposition}
\begin{proof}
First $G$ preserves the unsigned intersection pairing $Q$.
The module action and Proposition \ref{prop:pairing} give
$L_x^\dagger=L_x$ for the $Q$-adjoint, whereas $h^\dagger=-h$.
Taking adjoints of $[L_x,\Lambda_x]=h$ shows that
$\Lambda_x^\dagger$ is another degree-$-2$ lowering operator,
so $\Lambda_x^\dagger=\Lambda_x$. Hence all brackets
$[L_x,\Lambda_y]$ are $Q$-skew, and so is $\mathfrak s$.

Let $G_{\ZZ}$ be the stabiliser in $G(\QQ)$ of the chosen
lattice in $\IH$. It is arithmetic; its image under the
isogeny is commensurable with the integral orthogonal group
\cite[(5.4.5), (9.4.5)]{morris15}. After passing to a subgroup
$\Gamma_1$ of finite index, every action on $H^2$ has a lift
in $G_{\ZZ}$. Consider the group of pairs
\[
 \widetilde\Gamma_1=\{(\gamma,\sigma):\gamma\in\Gamma_1,
 \ \sigma\in G_{\ZZ},\ \bar\sigma=\gamma|_{H^2}\}.
\]
The operator $c(\gamma,\sigma)=\gamma\sigma^{-1}$ preserves
the lattice and $Q$ and commutes with all $L_x$. It commutes
with $h$ and therefore with the uniquely determined
$\Lambda_x$. It centralises the whole LLV algebra and $G$.
In particular $c$ is a homomorphism on the group of pairs.

Transport these operators to harmonic forms. By the star
formula after \eqref{eq:positive-pairing}, every $c$ commutes
with $*$. It preserves the positive inner product
$Q(\alpha,*\beta)$. The intersection of a compact orthogonal
group with a lattice stabiliser is finite. Thus the image of
$c$ is finite, and the image in $\Gamma_1$ of $\ker c$ has
finite index. Its action on $\IH$ is exactly the required
action through $G(\QQ)$.

For the last assertion take the Zariski closure of that
finite-index subgroup in $G_{\CC}$. Its projection contains
$\SO(W)$ because finite index preserves the identity component
of an algebraic closure. The isogeny has finite kernel; its
inverse image over $\SO(W)$ has connected component $G_W$.
A closed subgroup projecting onto $\SO(W)$ has the same Lie
algebra over this subgroup and contains $G_W$.
\end{proof}

Apply this proposition to locally trivial monodromy. Parallel
transport preserves the integral structures and the module action.
It also preserves the intersection pairing: on the regular parts
the complex orientations vary continuously in a locally trivial
family, so the oriented fundamental classes are a flat section.
Naturality of intersection duality then preserves $Q$. Thus
\begin{equation}\label{eq:monodromy-spin}
 \Mon'\subset G(\QQ),\qquad [\Mon:\Mon']<\infty,
 \qquad \overline{\Mon}^{\,\mathrm{Zar}}\supset G_{\CC}.
\end{equation}
The last inclusion uses the density on $H^2$ established in
Section \ref{sec:algebra}.

The proof is the harmonic-inner-product argument of
\cite[Theorem~3.5]{verbitsky13}; for moduli of $K3^{[n]}$ type
compare \cite[Lemma~4.13]{markman08}. It applies unchanged to
compact smooth hyperk\"ahler manifolds with their ordinary
cohomology and Poincar\'e pairing. This smooth version will be
used together with Proposition \ref{prop:spin} below.
